\documentclass[a4paper, 12pt]{article}
\usepackage[top=1in, bottom=1in, left=1in, right=1in]{geometry}
\usepackage{times}
\usepackage{setspace}
\onehalfspacing
\usepackage{amsmath, amssymb}
\usepackage{graphicx}
\usepackage{booktabs}
\usepackage{xcolor}
\usepackage{listings}
\usepackage{enumitem}
\usepackage[hidelinks]{hyperref}
\usepackage{fancyhdr}

\pagestyle{fancy}
\fancyhf{}
\fancyfoot[C]{\thepage}
\renewcommand{\headrulewidth}{0pt}

\definecolor{codebg}{RGB}{245,245,245}
\definecolor{codeblue}{RGB}{0,0,180}
\definecolor{codegreen}{RGB}{0,128,0}
\definecolor{codeorange}{RGB}{200,80,0}

\lstset{
  language=Python, backgroundcolor=\color{codebg},
  basicstyle=\small\ttfamily, breaklines=true,
  keywordstyle=\color{codeblue}\bfseries,
  commentstyle=\color{codegreen}\itshape,
  stringstyle=\color{codeorange},
  numbers=left, numberstyle=\tiny\color{gray},
  xleftmargin=2em, framexleftmargin=1.5em,
  frame=l, framerule=0.5pt, rulecolor=\color{gray!50},
}

\title{\textbf{Feature Extraction Pipeline: How Multi-Frame Video Features Are Extracted}\\[6pt]
\large A Technical Explanation for V-JEPA~2 and Qwen2.5-VL}
\author{Erize Rai}
\date{NOVEMBER, 2026}

\begin{document}
\maketitle

\section{Overview}

Both V-JEPA~2 and Qwen2.5-VL process \textbf{multiple frames per video clip}. However, each model produces a \textbf{single feature vector per clip} for the baseline classifiers. The key question is: \emph{how do multiple frames become one vector?}

The answer involves two stages:
\begin{enumerate}
    \item \textbf{Frame Sampling:} Multiple frames are uniformly sampled from each video clip.
    \item \textbf{Temporal Aggregation:} The model internally processes all frames \emph{jointly} as a spatiotemporal sequence, and the output is pooled (averaged) into a single vector.
\end{enumerate}

\section{Stage 1: Frame Sampling (Common to Both Models)}

For each video clip, frames are extracted using OpenCV's \texttt{VideoCapture}. Frames are sampled \textbf{uniformly} across the entire video duration (not consecutive):

\begin{lstlisting}
def load_frames_from_file(video_path, n_frames=8, size=224):
    cap = cv2.VideoCapture(video_path)
    total = int(cap.get(cv2.CAP_PROP_FRAME_COUNT))
    idxs = np.linspace(0, total - 1, n_frames, dtype=int)
    frames = []
    for i in idxs:
        cap.set(cv2.CAP_PROP_POS_FRAMES, i)
        ok, fr = cap.read()
        fr = cv2.cvtColor(fr, cv2.COLOR_BGR2RGB)
        fr = cv2.resize(fr, (size, size))
        frames.append(Image.fromarray(fr))
    return frames
\end{lstlisting}

\textbf{Key detail:} \texttt{np.linspace(0, total-1, n\_frames)} produces \textbf{evenly-spaced} frame indices. For example, a 120-frame video sampled at $n=8$ would select frames at indices [0, 17, 34, 51, 68, 85, 102, 119].

\begin{itemize}
    \item \textbf{V-JEPA~2:} 8 frames, resized to $224 \times 224$
    \item \textbf{Qwen2.5-VL:} 4 frames, resized to $224 \times 224$
\end{itemize}

\section{Stage 2: V-JEPA~2 Feature Extraction}

\subsection{How the Model Processes Multiple Frames}

V-JEPA~2 does \textbf{not} process each frame independently. The 8 frames are treated as a \textbf{3D video volume} and tokenized jointly:

\begin{enumerate}
    \item The 8 frames ($8 \times 224 \times 224 \times 3$) are passed to the \texttt{VJEPA2VideoProcessor}, which converts them into a single 4D pixel tensor.
    \item A \textbf{3D convolutional patch embedding} (\texttt{Conv3D}) with tubelet size $2 \times 16 \times 16$ converts the video into spatiotemporal tokens:
    \begin{itemize}
        \item Temporal dimension: $8 \div 2 = 4$ temporal tokens ($T'=4$)
        \item Spatial dimensions: $224 \div 16 = 14$ per axis, so $14 \times 14 = 196$ spatial tokens
        \item Total tokens: $4 \times 196 = 784$ tokens, each of dimension 1280
    \end{itemize}
    \item These 784 tokens pass through \textbf{all 32 transformer layers} of ViT-Huge as a single sequence --- the self-attention operates over \emph{all} spatiotemporal positions jointly, meaning \textbf{every frame attends to every other frame}.
\end{enumerate}

\subsection{How 784 Tokens Become One Vector (Baseline Classifiers)}

For the baseline classifiers (SVM, RF, LR, MLP, etc.), a single 1280-d vector per clip is needed. This is achieved via \textbf{global mean pooling over all 784 tokens} from the last hidden layer:

\begin{lstlisting}
# out.last_hidden_state shape: (1, 784, 1280)
single_feat = out.last_hidden_state.mean(dim=1)  # -> (1, 1280)
\end{lstlisting}

This averages across \textbf{all} 784 spatiotemporal tokens (all 4 temporal positions $\times$ all 196 spatial positions), producing a single 1280-dimensional feature vector that represents the entire 8-frame clip.

\subsection{How 784 Tokens Become Multi-Layer Features (HAF-STD Net)}

For the HAF-STD Net, the extraction preserves temporal structure:

\begin{enumerate}
    \item Hidden states are extracted from \textbf{4 intermediate layers} at 25\%, 50\%, 75\%, and 100\% depth (layers 8, 16, 24, 32 for ViT-Huge).
    \item At each layer, the 784 tokens are \textbf{reshaped} into a $4 \times 196$ grid (temporal $\times$ spatial).
    \item \textbf{Spatial mean pooling} averages across the 196 spatial positions per time step, yielding a $(T'=4, D=1280)$ matrix per layer.
    \item Stacking across 4 layers: final shape is $(4_\text{layers}, 4_\text{frames}, 1280)$ per clip.
\end{enumerate}

\begin{lstlisting}
# For each selected layer:
h = all_hidden[li].float()              # (1, 784, 1280)
h = h.view(1, 4, 196, 1280)            # reshape: (1, T', S, D)
h_temporal = h.mean(dim=2)             # spatial pool: (1, 4, 1280)
per_layer_frames.append(h_temporal)

# Stack all 4 layers:
stacked = torch.stack(per_layer_frames) # (4, 4, 1280)
\end{lstlisting}

\section{Stage 2: Qwen2.5-VL Feature Extraction}

\subsection{How the Model Processes Multiple Frames}

Qwen2.5-VL takes \textbf{4 frames} per clip and processes them via a vision-language pipeline:

\begin{enumerate}
    \item The 4 frames are formatted as a chat message with a minimal prompt (``Describe.'').
    \item The \textbf{Vision Encoder} converts each frame into visual tokens via $14 \times 14$ patches with a Conv3D temporal grouping (stride 2), producing visual tokens.
    \item These visual tokens are projected into the LM embedding space (3584-d) and \textbf{interleaved with text tokens} (the prompt).
    \item The full multimodal sequence passes through the \textbf{Qwen2.5 LM Decoder} (28 transformer layers).
\end{enumerate}

\subsection{How the Token Sequence Becomes One Vector}

The hidden states from the \textbf{final decoder layer} are mean-pooled across the \emph{entire} sequence dimension (both visual and text tokens):

\begin{lstlisting}
# out.hidden_states[-1] shape: (1, seq_len, 3584)
feat = out.hidden_states[-1].mean(dim=1)  # -> (1, 3584)
\end{lstlisting}

This produces a single \textbf{3584-dimensional} feature vector per clip.

\section{Summary: The Complete Pipeline}

\begin{center}
\renewcommand{\arraystretch}{1.4}
\begin{tabular}{p{4cm} p{5cm} p{5cm}}
\toprule
\textbf{Stage} & \textbf{V-JEPA~2} & \textbf{Qwen2.5-VL} \\
\midrule
Input frames & 8 frames at $224 \times 224$ & 4 frames at $224 \times 224$ \\
\midrule
Tokenization & Conv3D tubelets ($2\!\times\!16\!\times\!16$) $\rightarrow$ 784 tokens & Conv3D patches ($14\!\times\!14$, stride 2) $\rightarrow$ variable tokens \\
\midrule
Encoder & 32-layer ViT-Huge (654M params) & Vision Encoder + 28-layer LM Decoder (7B params, 4-bit) \\
\midrule
Temporal handling & All 8 frames jointly via 3D spatiotemporal attention & Frames interleaved with text in autoregressive decoder \\
\midrule
Baseline feature & Mean pool over all 784 tokens $\rightarrow$ \textbf{1280-d} & Mean pool over all sequence tokens $\rightarrow$ \textbf{3584-d} \\
\midrule
HAF-STD feature & 4 layers $\times$ spatial pool $\rightarrow$ \textbf{(4, 4, 1280)} & Not used (V-JEPA~2 only) \\
\bottomrule
\end{tabular}
\end{center}

\section{Why Only 8 Frames from Videos with Hundreds of Frames?}

A natural question arises: if a video clip has 100--1000+ frames (at 24--30 fps), why sample only 8? There are four complementary reasons:

\subsection{The Model Architecture Requires It}

V-JEPA~2 was \textbf{pretrained} with 8-frame input clips. Its Conv3D tubelet embedding layer ($2 \times 16 \times 16$) is architecturally designed for this temporal depth. The tubelet's temporal stride of 2 over 8 frames produces exactly $T' = 4$ temporal tokens --- this is a fixed architectural parameter, not an arbitrary choice. Using a different number of frames would produce a mismatched token sequence that does not align with the pretrained positional embeddings and attention patterns.

Similarly, Qwen2.5-VL's video processing pipeline uses 4 frames by default, matching its Conv3D temporal grouping stride.

\subsection{Uniform Sampling Covers the Full Temporal Span}

The 8 frames are \textbf{not consecutive} --- they are \textbf{uniformly spaced} across the entire video duration:

\begin{lstlisting}
idxs = np.linspace(0, total - 1, n_frames, dtype=int)
\end{lstlisting}

For a 240-frame video (8 seconds at 30 fps), the 8 selected frame indices would be:
$$[0,\; 34,\; 68,\; 103,\; 137,\; 171,\; 205,\; 239]$$

This means the model sees the \textbf{beginning, middle, and end} of the entire clip. Each sampled frame is separated by approximately 1 second, covering the full temporal extent of the action.

\begin{center}
\small
\setlength{\unitlength}{0.9pt}
\begin{picture}(380,50)
\put(0,25){\line(1,0){380}}
\put(0,20){\line(0,1){10}}
\put(380,20){\line(0,1){10}}
\put(0,10){\makebox(0,0){\scriptsize 0}}
\put(380,10){\makebox(0,0){\scriptsize 239}}
% sampled frames
\put(0,25){\circle*{6}}
\put(54,25){\circle*{6}}
\put(108,25){\circle*{6}}
\put(163,25){\circle*{6}}
\put(217,25){\circle*{6}}
\put(271,25){\circle*{6}}
\put(325,25){\circle*{6}}
\put(380,25){\circle*{6}}
% labels
\put(0,40){\makebox(0,0){\scriptsize f1}}
\put(54,40){\makebox(0,0){\scriptsize f2}}
\put(108,40){\makebox(0,0){\scriptsize f3}}
\put(163,40){\makebox(0,0){\scriptsize f4}}
\put(217,40){\makebox(0,0){\scriptsize f5}}
\put(271,40){\makebox(0,0){\scriptsize f6}}
\put(325,40){\makebox(0,0){\scriptsize f7}}
\put(380,40){\makebox(0,0){\scriptsize f8}}
\end{picture}

\textit{8 uniformly-sampled frames (black dots) from a 240-frame video timeline.}
\end{center}

\subsection{GPU Memory Constraints}

Self-attention has \textbf{quadratic complexity} $\mathcal{O}(N^2)$ in the number of tokens. With 8 frames, V-JEPA~2 produces 784 tokens. Doubling to 16 frames would produce 1,568 tokens, requiring approximately $4\times$ more memory for attention computation --- infeasible on the 15.6~GB Tesla T4 GPU used in this study, especially alongside the 4-bit quantized 7B-parameter Qwen2.5-VL model.

\begin{center}
\renewcommand{\arraystretch}{1.3}
\begin{tabular}{ccc}
\toprule
\textbf{Frames} & \textbf{Tokens} & \textbf{Relative Attention Cost} \\
\midrule
4  & 392  & $0.25\times$ \\
8  & 784  & $1.0\times$ (used) \\
16 & 1568 & $4.0\times$ \\
32 & 3136 & $16.0\times$ \\
\bottomrule
\end{tabular}
\end{center}

\subsection{Temporal Redundancy in Video}

At 24--30 fps, adjacent frames differ by only 33--42 milliseconds --- the visual difference between consecutive frames is \textbf{negligible}. For action recognition tasks like distinguishing Balance Beam from Vault, the discriminative information lies in \emph{key body poses and transitions}, not in every sub-frame motion. Eight uniformly-spaced frames capture the essential ``stroboscopic'' motion trajectory.

\subsection{This Is the Field-Wide Standard}

Virtually all modern video transformers use 8--16 sampled frames:

\begin{center}
\renewcommand{\arraystretch}{1.3}
\begin{tabular}{lc}
\toprule
\textbf{Model} & \textbf{Default Frames} \\
\midrule
V-JEPA 2 / V-JEPA 2.1 (Meta) & 8 \\
VideoMAE (Tong et al., 2022)  & 16 \\
TimeSformer (Bertasius et al., 2021) & 8 \\
ViViT (Arnab et al., 2021) & 8--32 \\
Qwen2.5-VL (Alibaba) & 4 \\
\bottomrule
\end{tabular}
\end{center}

This is not a limitation of the present study --- it is the established methodology for frozen-feature video representation evaluation.

\section{Key Takeaway}

\textbf{Multiple frames are NOT processed independently.} Both models treat multi-frame input as a \emph{joint spatiotemporal sequence}. The self-attention in V-JEPA~2's ViT-Huge operates over all 784 tokens from all 8 frames simultaneously --- frame 1's patches attend to frame 8's patches. The final feature vector is obtained by \textbf{mean-pooling} this joint representation, which inherently captures cross-frame temporal relationships learned during pretraining. This is fundamentally different from extracting a per-frame feature and averaging afterward.

The choice of 8 frames is architecturally determined by the pretrained model, covers the full temporal span via uniform sampling, and is the standard methodology across the video understanding research community.

\end{document}
